To extend the construction to an arbitrary spin, consider once more a system
of $2s$ particles of spin $1/2$, and express its wave function as a product of
$2s$ spinors. The spin operator for this system is the sum of the individual
particles' spin operators. Each term acts only on its corresponding spinor,
with the result specified by formulas (57.4). Returning then to arbitrary
symmetric spinors, that is, to the wave functions of a particle of spin $s$,
gives
\begin{equation}
 \begin{aligned}
 (\hat s_x\psi)^{\overbrace{11\ldots}^{s+\sigma}
                         \overbrace{22\ldots}^{s-\sigma}}
 &=\frac{s+\sigma}{2}\psi^{\overbrace{11\ldots}^{s+\sigma-1}
                                  \overbrace{22\ldots}^{s-\sigma+1}}
 +\frac{s-\sigma}{2}\psi^{\overbrace{11\ldots}^{s+\sigma+1}
                                  \overbrace{22\ldots}^{s-\sigma-1}},\\
 (\hat s_y\psi)^{\overbrace{11\ldots}^{s+\sigma}
                         \overbrace{22\ldots}^{s-\sigma}}
 &=-i\frac{s+\sigma}{2}\psi^{\overbrace{11\ldots}^{s+\sigma-1}
                                   \overbrace{22\ldots}^{s-\sigma+1}}
 +i\frac{s-\sigma}{2}\psi^{\overbrace{11\ldots}^{s+\sigma+1}
                                  \overbrace{22\ldots}^{s-\sigma-1}},\\
 (\hat s_z\psi)^{\overbrace{11\ldots}^{s+\sigma}
                         \overbrace{22\ldots}^{s-\sigma}}
 &=\sigma\psi^{\overbrace{11\ldots}^{s+\sigma}
                       \overbrace{22\ldots}^{s-\sigma}}.
 \end{aligned}
 \tag{57.5}
\end{equation}
